% !TEX root = ../main.tex
%% Chapitre 14 — Groupe symétrique et déterminants
%% Programme officiel MPSI/MP2I (BO spécial n°1 du 11/02/2021), deuxième semestre :
%% A) groupe symétrique (cycles, transpositions, décomposition, signature) ;
%% B) déterminants (formes n-linéaires alternées, déterminant dans une base,
%% d'un endomorphisme, d'une matrice ; calculs : opérations, cofacteurs,
%% triangulaires, Vandermonde ; comatrice).
%% Résultats vérifiés par verifications/chapitre_14.py.

\begin{ouverture}[determinants]{14}{Groupe symétrique et déterminants}{Partie 2 · Deuxième semestre}
  {Permutations\enspace·\enspace Cycles\enspace·\enspace Signature\enspace·\enspace
   Formes alternées\enspace·\enspace Déterminants\enspace·\enspace Comatrice}
  \objectif{Décomposer une permutation en cycles à supports disjoints, puis en transpositions.}
  \objectif{Calculer une signature et utiliser le morphisme $\varepsilon$.}
  \objectif{Comprendre le déterminant comme forme $n$-linéaire alternée — une aire, un volume.}
  \objectif{Calculer un déterminant : opérations élémentaires, développement, triangulaires.}
  \objectif{Caractériser bases, automorphismes et matrices inversibles par le déterminant.}
  \objectif{Connaître Vandermonde et la comatrice, et savoir s'en servir.}
  \prerequis{Groupes et morphismes (chapitre 8) ; le groupe symétrique prolonge les bijections (chapitre 1).}
  \prerequis{Matrices : produit, inversibilité, rang, similitude (chapitres 9 et 13).}
  \prerequis{Espaces vectoriels : bases, dimension (chapitre 12).}
\end{ouverture}

%% =====================================================================
%%  L'ESSENTIEL
%% =====================================================================
\essentiel

\essSection{1}{Le groupe symétrique $S_n$}

\begin{definition}[Permutations, cycles, transpositions]
$S_n$ est le groupe des bijections de $\ens{1,\dots,n}$, pour la
composition ; son cardinal est $n!$ et il n'est pas commutatif dès que
$n\geqslant3$. Un \textbf{cycle} $(a_1\ a_2\ \dots\ a_p)$ envoie
$a_1\mapsto a_2\mapsto\dots\mapsto a_p\mapsto a_1$ et fixe le reste ; son
\emph{support} est $\ens{a_1,\dots,a_p}$. Une \textbf{transposition} est un
cycle de longueur $2$ : un simple échange.
\end{definition}

\begin{theoreme}[Décomposition en cycles]
Toute permutation se décompose en produit de cycles à supports
\emph{disjoints}, de façon unique à l'ordre près — et des cycles à supports
disjoints commutent (démonstration non exigible, mais la décomposition doit
être un réflexe : suivre les orbites de chaque élément). Toute permutation
est aussi un produit de transpositions : il suffit de décomposer chaque
cycle, $(a_1\ \dots\ a_p)=(a_1\ a_p)\cdots(a_1\ a_3)(a_1\ a_2)$.
\end{theoreme}

\begin{exemple}
Dans $S_6$, $\sigma=\begin{pmatrix}1&2&3&4&5&6\\3&5&4&1&2&6\end{pmatrix}$ :
l'orbite de $1$ est $1\to3\to4\to1$, celle de $2$ est $2\to5\to2$, et $6$
est fixe : $\sigma=(1\ 3\ 4)(2\ 5)$. Son ordre est
$\operatorname{ppcm}(3,2)=6$ (chapitres 7 et 8).
\end{exemple}

\essSection{2}{La signature}

\begin{theoreme}[Signature]
Il existe un \emph{unique} morphisme de groupes
$\varepsilon:S_n\to\ens{-1,1}$ envoyant toute transposition sur $-1$
(démonstration non exigible) : la \textbf{signature}. Pour un cycle de
longueur $p$ : $\varepsilon=(-1)^{p-1}$ ; pour un produit, on multiplie.
Une permutation est \emph{paire} si $\varepsilon(\sigma)=1$ : le nombre de
transpositions d'une décomposition a toujours la même parité.
\end{theoreme}

\begin{exemple}
$\varepsilon\bigl((1\ 3\ 4)(2\ 5)\bigr)=(+1)\times(-1)=-1$ ;
$\varepsilon\bigl((1\ 2\ 3\ 4\ 5)(6\ 7)\bigr)=(+1)(-1)=-1$. On peut aussi
compter les \emph{inversions} : $\varepsilon(\sigma)=(-1)^{N(\sigma)}$ où
$N(\sigma)$ est le nombre de paires $i<j$ avec $\sigma(i)>\sigma(j)$
(exercice 17).
\end{exemple}

\essSection{3}{Formes $n$-linéaires alternées}

\begin{definition}[Forme $n$-linéaire alternée]
Sur un $\K$-espace $E$ de dimension $n$, une forme $n$-linéaire
$f(x_1,\dots,x_n)$ est \textbf{alternée} si elle s'annule dès que deux
arguments sont égaux. Elle est alors \emph{antisymétrique} (échanger deux
arguments change le signe), et pour toute permutation :
$f(x_{\sigma(1)},\dots,x_{\sigma(n)})=\varepsilon(\sigma)\,f(x_1,\dots,x_n)$.
Si la famille $(x_1,\dots,x_n)$ est \emph{liée}, alors
$f(x_1,\dots,x_n)=0$.
\end{definition}

\noindent
\begin{minipage}[c]{0.55\linewidth}
\begin{theoreme}[Déterminant dans une base]
Soit $e$ une base de $E$. Il existe une \emph{unique} forme $n$-linéaire
alternée $\det_e$ telle que $\det_e(e)=1$ (existence non exigible), et
toute forme $n$-linéaire alternée lui est proportionnelle. En
coordonnées :
\[
\det_e(x_1,\dots,x_n)
=\sum_{\sigma\in S_n}\varepsilon(\sigma)\prod_{i=1}^{n}a_{\sigma(i),i},
\]
où $a_{ij}$ est la $i$-ième coordonnée de $x_j$. Pour deux bases :
$\det_{e'}=\det_{e'}(e)\det_e$, et
\[
(x_1,\dots,x_n)\ \text{base}\iff\det_e(x_1,\dots,x_n)\neq0 .
\]
\end{theoreme}
\end{minipage}\hfill
\begin{minipage}[c]{0.41\linewidth}
\centering
\begin{tikzpicture}[scale=0.8,line cap=round]
  \draw[otMuted,line width=0.5pt,-{Stealth[length=1.6mm]}] (-0.4,0) -- (3.9,0);
  \draw[otMuted,line width=0.5pt,-{Stealth[length=1.6mm]}] (0,-0.4) -- (0,3.3);
  \fill[otGreen,opacity=0.3] (0,0) -- (3,1) -- (4,3) -- (1,2) -- cycle;
  \draw[otGreenDark,line width=0.8pt,line join=round] (0,0) -- (3,1) -- (4,3) -- (1,2) -- cycle;
  \draw[otNavy,line width=1.1pt,-{Stealth[length=2mm]}] (0,0) -- (3,1);
  \draw[otNavy,line width=1.1pt,-{Stealth[length=2mm]}] (0,0) -- (1,2);
  \node[font=\scriptsize,text=otNavy] at (3.15,0.7) {$u$};
  \node[font=\scriptsize,text=otNavy] at (0.68,2.05) {$v$};
  \node[font=\scriptsize,text=otGreenDark] at (2,1.5) {$\det(u,v)$};
\end{tikzpicture}\\[0.5mm]
\legende{Dans $\R^{2}$ : $\det(u,v)$ est l'\emph{aire orientée} du
parallélogramme ; dans $\R^{3}$, le volume orienté du parallélépipède.}
\end{minipage}

\essSection{4}{Déterminant d'un endomorphisme, d'une matrice}

\begin{theoreme}[Déterminant d'un endomorphisme]
Pour $u\in\mathcal{L}(E)$, le scalaire
$\det u=\det_e\bigl(u(e_1),\dots,u(e_n)\bigr)$ ne dépend pas de la base
$e$ : c'est le facteur par lequel $u$ multiplie les volumes orientés. On a
$\det(v\circ u)=\det v\cdot\det u$, $\det(\mathrm{id})=1$, et
\[
u\ \text{automorphisme}\iff\det u\neq0 ,
\qquad\text{avec alors}\quad\det\bigl(u^{-1}\bigr)=\frac{1}{\det u} .
\]
\end{theoreme}

\begin{theoreme}[Déterminant d'une matrice carrée]
$\det A$ est le déterminant des colonnes de $A$ dans la base canonique de
$\K^{n}$ : il est $n$-linéaire alterné en les \emph{colonnes}, et aussi en
les \emph{lignes} car $\det\bigl(A^{\mathsf T}\bigr)=\det A$. Propriétés :
\[
\det(AB)=\det A\det B,\qquad
\det(\lambda A)=\lambda^{n}\det A,\qquad
A\in\operatorname{GL}_n(\K)\iff\det A\neq0 ,
\]
et $\det:\operatorname{GL}_n(\K)\to\K^{*}$ est un morphisme de groupes.
Deux matrices semblables ont le même déterminant : avec le rang et la
trace, c'est le troisième invariant de similitude (chapitre 13).
\end{theoreme}

\essSection{5}{Les calculs en pratique}

\begin{propriete}[Opérations élémentaires et développement]
Échanger deux colonnes (ou lignes) multiplie le déterminant par $-1$ ;
multiplier une colonne par $\lambda$ le multiplie par $\lambda$ ; ajouter à
une colonne un multiple d'une autre \emph{ne change rien} — l'arme
principale. Le \textbf{cofacteur} d'indice $(i,j)$ est
$A_{ij}=(-1)^{i+j}\Delta_{ij}$, où $\Delta_{ij}$ est le déterminant extrait
en rayant ligne $i$ et colonne $j$, et l'on peut \emph{développer} par
rapport à toute ligne ou colonne :
$\det A=\sum_{i}a_{ij}A_{ij}$ (colonne $j$ fixée). Enfin, le déterminant
d'une matrice \emph{triangulaire} est le produit de ses coefficients
diagonaux.
\end{propriete}

\begin{theoreme}[Déterminant de Vandermonde]
Pour $a_1,\dots,a_n\in\K$ :
\[
\begin{vmatrix}
1&1&\cdots&1\\
a_1&a_2&\cdots&a_n\\
\vdots&\vdots&&\vdots\\
a_1^{\,n-1}&a_2^{\,n-1}&\cdots&a_n^{\,n-1}
\end{vmatrix}
=\prod_{1\leqslant i<j\leqslant n}(a_j-a_i) :
\]
non nul si et seulement si les $a_i$ sont deux à deux distincts. C'est
l'inversibilité du système d'interpolation : l'existence et l'unicité du
polynôme interpolateur de Lagrange (chapitres 10 et 12) en découlent une
troisième fois.
\end{theoreme}

\begin{savoirfaire}[Calculer un déterminant sans souffrir]
D'abord \emph{simplifier} : factoriser une ligne, créer des zéros par
opérations élémentaires (souvent $C_j\leftarrow C_j-C_1$ ou sommer toutes
les colonnes quand la somme est constante), viser une forme triangulaire.
Ne développer par rapport à une ligne ou colonne que lorsqu'elle est déjà
presque vide. Pour un déterminant à paramètre, factoriser le résultat
avant de discuter.
\end{savoirfaire}

\essSection{6}{La comatrice}

\begin{theoreme}[Comatrice et inverse]
La \textbf{comatrice} $\operatorname{Com}A$ est la matrice des cofacteurs
$\bigl(A_{ij}\bigr)$. Pour toute $A\in\mathcal{M}_n(\K)$ :
\[
A\,(\operatorname{Com}A)^{\mathsf T}
=(\operatorname{Com}A)^{\mathsf T}A=\det(A)\,I_n ,
\qquad\text{d'où, si }\det A\neq0 :\quad
A^{-1}=\frac{1}{\det A}\,(\operatorname{Com}A)^{\mathsf T}.
\]
En taille $2$ : $\begin{pmatrix}a&b\\c&d\end{pmatrix}^{-1}
=\dfrac{1}{ad-bc}\begin{pmatrix}d&-b\\-c&a\end{pmatrix}$ — à connaître par
cœur. Au-delà de $n=3$, la formule sert surtout à la théorie (exercice 22) ;
pour le calcul effectif, le pivot reste roi.
\end{theoreme}

%% =====================================================================
%%  EXERCICES
%% =====================================================================
\exercices
\rubrique{Applications}

\begin{exercice}[1]{Composer des permutations}
Dans $S_4$, soient
$\sigma=\begin{pmatrix}1&2&3&4\\2&3&1&4\end{pmatrix}$ et
$\tau=\begin{pmatrix}1&2&3&4\\2&1&4&3\end{pmatrix}$.
\begin{questions}
  \item Calculer $\sigma\circ\tau$ et $\tau\circ\sigma$. Commutent-elles ?
  \item Déterminer $\sigma^{-1}$.
  \item Écrire $\sigma$ et $\tau$ comme produits de cycles.
\end{questions}
\end{exercice}

\begin{exercice}[1]{Décomposer en cycles}
Décomposer en cycles à supports disjoints, puis donner l'ordre :
\begin{questions}
  \item $\sigma_1=\begin{pmatrix}1&2&3&4&5&6\\3&5&4&1&2&6\end{pmatrix}$ ;
  \item $\sigma_2=\begin{pmatrix}1&2&3&4&5&6&7\\2&3&4&5&1&7&6\end{pmatrix}$.
\end{questions}
\end{exercice}

\begin{exercice}[1]{Signatures}
\begin{questions}
  \item Rappeler la signature d'un cycle de longueur $p$, et la retrouver
    en le décomposant en transpositions.
  \item Calculer $\varepsilon(\sigma_1)$ et $\varepsilon(\sigma_2)$
    (exercice 2).
  \item Vérifier sur $\sigma$ et $\tau$ de l'exercice 1 que
    $\varepsilon(\sigma\tau)=\varepsilon(\sigma)\varepsilon(\tau)$.
\end{questions}
\end{exercice}

\begin{exercice}[1]{Déterminants 2 × 2 et 3 × 3}
\begin{questions}
  \item Calculer $\begin{vmatrix}2&3\\1&4\end{vmatrix}$ et
    $\begin{vmatrix}1&2&3\\4&5&6\\7&8&10\end{vmatrix}$.
  \item Vérifier sur $A=\begin{pmatrix}1&2&3\\4&5&6\\7&8&10\end{pmatrix}$
    et $B=\begin{pmatrix}1&0&2\\0&1&1\\1&1&0\end{pmatrix}$ que
    $\det(AB)=\det A\det B$.
\end{questions}
\end{exercice}

\begin{exercice}[1]{Créer des zéros}
Calculer par opérations élémentaires :
\begin{questions}
  \item $\begin{vmatrix}1&1&1\\1&2&3\\1&3&6\end{vmatrix}$ ;
  \item $\begin{vmatrix}0&1&2\\1&2&3\\2&3&4\end{vmatrix}$. Pourquoi ce
    résultat était-il prévisible ?
\end{questions}
\end{exercice}

\begin{exercice}[1]{Triangulaires, transposée, inverse}
Soit $T=\begin{pmatrix}2&1&5\\0&3&-1\\0&0&4\end{pmatrix}$.
\begin{questions}
  \item Donner $\det T$, puis $\det\bigl(T^{\mathsf T}\bigr)$ et
    $\det\bigl(T^{-1}\bigr)$.
  \item Que vaut $\det(2T)$ ? (Attention au piège.)
\end{questions}
\end{exercice}

\begin{exercice}[1]{Développer intelligemment}
Calculer
$\begin{vmatrix}2&0&0&1\\0&3&0&0\\0&0&1&4\\1&0&0&1\end{vmatrix}$
en développant par rapport à une ligne ou colonne bien choisie.
\end{exercice}

\begin{exercice}[1]{Aires et volumes orientés}
\begin{questions}
  \item Calculer l'aire du parallélogramme construit sur $u=(3,1)$ et
    $v=(1,2)$, puis celle du triangle $(0,u,u+v)$.
  \item Calculer le volume du parallélépipède construit sur
    $(1,0,0)$, $(1,1,0)$ et $(1,1,1)$.
  \item Que signifie le signe de $\det(u,v)$ ? Donner un exemple de base
    indirecte de $\R^{2}$.
\end{questions}
\end{exercice}

\rubrique{Entraînement}

\begin{exercice}[2]{La formule de Leibniz à la main}
\begin{questions}
  \item Écrire la formule de Leibniz pour $n=2$ et retrouver $ad-bc$.
  \item Pour $n=3$, expliciter les $6$ termes et retrouver la règle de
    Sarrus ; vérifier sur le déterminant $3\times3$ de l'exercice 4.
  \item Combien la formule compte-t-elle de termes pour $n=10$ ?
    Commenter l'intérêt pratique.
\end{questions}
\end{exercice}

\begin{exercice}[2]{Base ou pas ? Le critère du déterminant}
\begin{questions}
  \item La famille $\bigl((1,0,1),(1,1,0),(0,1,1)\bigr)$ est-elle une base
    de $\R^{3}$ ? Comparer avec la méthode du chapitre 12.
  \item Même question pour $(1+X,\,1-X,\,X^{2})$ dans $\R_2[X]$, en
    passant par les coordonnées.
\end{questions}
\end{exercice}

\begin{exercice}[2]{Déterminants à paramètre}
Pour $t\in\R$, calculer et factoriser :
\begin{questions}
  \item $\begin{vmatrix}1&t\\t&1\end{vmatrix}$ ;
  \item $\begin{vmatrix}1&1&t\\1&t&1\\t&1&1\end{vmatrix}$ — retrouver la
    discussion de rang du chapitre 13.
\end{questions}
\end{exercice}

\begin{exercice}[2]{Le déterminant de Vandermonde en taille 3}
\begin{questions}
  \item Montrer par opérations élémentaires que
    $\begin{vmatrix}1&1&1\\a&b&c\\a^{2}&b^{2}&c^{2}\end{vmatrix}
    =(b-a)(c-a)(c-b)$.
  \item En déduire que trois réels distincts donnent toujours un système
    d'interpolation inversible, et faire le lien avec Lagrange
    (chapitres 10 et 12).
\end{questions}
\end{exercice}

\begin{exercice}[2]{Matrices de permutation}
Pour $\sigma\in S_n$, soit $P_\sigma$ la matrice dont la colonne $j$ est
$e_{\sigma(j)}$.
\begin{questions}
  \item Vérifier que $P_\sigma P_\tau=P_{\sigma\tau}$ et que
    $P_\sigma^{-1}=P_\sigma^{\mathsf T}$.
  \item Montrer que $\det P_\sigma=\varepsilon(\sigma)$ : la signature
    est un déterminant.
  \item Que compte $\operatorname{tr}P_\sigma$ ?
\end{questions}
\end{exercice}

\begin{exercice}[2]{Comatrice et inversion}
Soit $A=\begin{pmatrix}1&2&0\\0&1&1\\1&0&1\end{pmatrix}$.
\begin{questions}
  \item Calculer $\det A$, puis $\operatorname{Com}A$.
  \item Vérifier $A\,(\operatorname{Com}A)^{\mathsf T}=\det(A)\,I_3$ et en
    déduire $A^{-1}$.
  \item Retrouver la formule d'inversion des matrices $2\times2$.
\end{questions}
\end{exercice}

\begin{exercice}[2]{Déterminants d'endomorphismes}
Calculer le déterminant des endomorphismes suivants de $\R_2[X]$, et dire
lesquels sont des automorphismes :
\begin{questions}
  \item $u:P\mapsto P+P'$ ;
  \item $v:P\mapsto P(X+1)$ ;
  \item $w:P\mapsto XP'$.
\end{questions}
\end{exercice}

\begin{exercice}[2]{Un déterminant circulant}
Soit $j=\mathrm{e}^{2\mathrm{i}\pi/3}$ et $a,b,c\in\Cx$.
\begin{questions}
  \item Montrer que
    $\begin{vmatrix}a&b&c\\c&a&b\\b&c&a\end{vmatrix}
    =a^{3}+b^{3}+c^{3}-3abc$.
  \item En remarquant que les colonnes $(1,1,1)$, $(1,j,j^{2})$ et
    $(1,j^{2},j)$ sont des vecteurs « propres » de la permutation
    circulaire, établir la factorisation
    $a^{3}+b^{3}+c^{3}-3abc=(a+b+c)(a+jb+j^{2}c)(a+j^{2}b+jc)$
    (chapitre 3).
\end{questions}
\end{exercice}

\begin{exercice}[2]{La signature par les inversions}
Pour $\sigma\in S_n$, on note $N(\sigma)$ le nombre d'\emph{inversions} :
les paires $i<j$ telles que $\sigma(i)>\sigma(j)$.
\begin{questions}
  \item Calculer $N(\sigma)$ pour
    $\sigma=\begin{pmatrix}1&2&3\\3&1&2\end{pmatrix}$, puis pour une
    transposition de $S_2$.
  \item Montrer que composer $\sigma$ par une transposition de voisins
    $(k\ k{+}1)$ change $N$ de $\pm1$.
  \item En déduire $\varepsilon(\sigma)=(-1)^{N(\sigma)}$.
\end{questions}
\end{exercice}

\begin{exercice}[2]{Un déterminant tridiagonal}
Pour $n\geqslant1$, soit $D_n$ le déterminant de la matrice $n\times n$
portant des $2$ sur la diagonale et des $1$ sur les deux sous-diagonales
adjacentes.
\begin{questions}
  \item Calculer $D_1$ et $D_2$, puis montrer en développant que
    $D_n=2D_{n-1}-D_{n-2}$.
  \item En déduire $D_n=n+1$ (chapitre 5 : récurrences linéaires d'ordre
    2).
\end{questions}
\end{exercice}

\rubrique{Approfondissement}

\begin{exercice}[3]{Générateurs de $S_n$}
\begin{questions}
  \item Montrer que les transpositions $(1\ 2),(1\ 3),\dots,(1\ n)$
    engendrent $S_n$.
  \item Montrer que $(1\ 2)$ et le cycle $c=(1\ 2\ \dots\ n)$ suffisent.
    \emph{Calculer $c^{k}(1\ 2)c^{-k}$.}
\end{questions}
\end{exercice}

\begin{exercice}[3]{L'ordre maximal dans $S_n$}
\begin{questions}
  \item Montrer que l'ordre d'une permutation est le ppcm des longueurs de
    ses cycles (chapitres 7 et 8).
  \item Déterminer l'ordre maximal d'un élément de $S_5$, puis de $S_7$,
    en passant en revue les partitions de $5$ et de $7$.
\end{questions}
\end{exercice}

\begin{exercice}[3]{$\det\bigl(A^{\mathsf T}\bigr)=\det A$, par Leibniz}
\begin{questions}
  \item Dans la formule de Leibniz appliquée à $A^{\mathsf T}$, effectuer
    le changement d'indice $\sigma\mapsto\sigma^{-1}$ et conclure.
    \emph{On utilisera $\varepsilon(\sigma^{-1})=\varepsilon(\sigma)$.}
  \item En déduire que le déterminant est aussi alterné en les lignes, et
    que toute propriété « colonnes » se transpose aux lignes.
\end{questions}
\end{exercice}

\begin{exercice}[3]{Inversibles dans $\mathcal{M}_n(\Z)$}
Soit $A$ une matrice carrée à coefficients \emph{entiers}.
\begin{questions}
  \item Montrer que $A$ possède une inverse à coefficients entiers si et
    seulement si $\det A=\pm1$. \emph{Penser à la comatrice.}
  \item Vérifier sur $\begin{pmatrix}2&1\\1&1\end{pmatrix}$, puis trouver
    une matrice entière inversible dans $\mathcal{M}_2(\R)$ mais pas dans
    $\mathcal{M}_2(\Z)$.
  \item Quel lien avec le théorème de Bézout (chapitre 7) en taille 1 ?
\end{questions}
\end{exercice}

\begin{exercice}[3]{Le déterminant transforme les aires}
On admet qu'une application linéaire $u$ de $\R^{2}$ multiplie les aires
par $|\det u|$ (c'est la définition géométrique du déterminant).
\begin{questions}
  \item Retrouver ainsi l'aire du parallélogramme de l'exercice 8.
  \item L'ellipse $\dfrac{x^{2}}{a^{2}}+\dfrac{y^{2}}{b^{2}}\leqslant1$ est
    l'image du disque unité par $\operatorname{diag}(a,b)$ : en déduire que
    son aire vaut $\pi ab$.
  \item Pourquoi $\det(AB)=\det A\det B$ devient-il évident de ce point de
    vue ?
\end{questions}
\end{exercice}

\begin{exercice}[3]{Vandermonde en toute taille}
On note $V(a_1,\dots,a_n)$ le déterminant de Vandermonde.
\begin{questions}
  \item Fixant $a_1,\dots,a_{n-1}$ distincts, montrer que
    $x\mapsto V(a_1,\dots,a_{n-1},x)$ est une fonction polynomiale de degré
    $n-1$, de racines $a_1,\dots,a_{n-1}$, et de coefficient dominant
    $V(a_1,\dots,a_{n-1})$.
  \item En déduire par récurrence la formule
    $V(a_1,\dots,a_n)=\prod_{i<j}(a_j-a_i)$.
  \item Vérifier la formule pour $(0,1,2,3)$.
\end{questions}
\end{exercice}

% !TEX root = ../../main.tex
%% Défi du chapitre 14 — Le déterminant de Cauchy et les matrices de Hilbert
%% Résultats vérifiés par verifications/defis_11_15.py.

\rubrique{Défi}

\begin{defi}{Le déterminant de Cauchy et les matrices de Hilbert}{2 h}
Soient $a_1,\dots,a_n$ et $b_1,\dots,b_n$ des réels tels que
$a_i+b_j\neq0$ pour tous $i,j$. On note $C_n$ le déterminant de la matrice
$\Bigl(\dfrac{1}{a_i+b_j}\Bigr)_{1\leqslant i,j\leqslant n}$, appelé
\emph{déterminant de Cauchy}.

\partieExo{A}{Petites tailles}
\begin{questions}
  \item Calculer $C_1$, puis $C_2$, et vérifier que
    \[
    C_2=\frac{(a_2-a_1)(b_2-b_1)}{(a_1+b_1)(a_1+b_2)(a_2+b_1)(a_2+b_2)} .
    \]
\end{questions}

\partieExo{B}{La formule de Cauchy}
\begin{questions}[resume]
  \item Pour $j<n$, on retranche la dernière colonne à la colonne $j$.
    Montrer que le nouveau coefficient $(i,j)$ vaut
    $\dfrac{b_n-b_j}{(a_i+b_j)(a_i+b_n)}$.
  \item Factoriser $b_n-b_j$ dans chaque colonne $j<n$ et
    $\frac{1}{a_i+b_n}$ dans chaque ligne $i$ : la dernière colonne ne
    contient plus que des $1$. Retrancher alors la dernière ligne à chacune
    des autres, factoriser de nouveau, et développer suivant la dernière
    colonne pour obtenir
    \[
    C_n=\frac{\prod_{k=1}^{n-1}(a_n-a_k)(b_n-b_k)}
    {\prod_{k=1}^{n}(a_k+b_n)\ \prod_{k=1}^{n-1}(a_n+b_k)}\;C_{n-1} .
    \]
  \item En déduire la \emph{formule de Cauchy} :
    \[
    C_n=\frac{\prod_{1\leqslant i<j\leqslant n}(a_j-a_i)(b_j-b_i)}
    {\prod_{1\leqslant i,j\leqslant n}(a_i+b_j)} .
    \]
    Quand $C_n$ est-il non nul ?
\end{questions}

\partieExo{C}{Les matrices de Hilbert}
On note $H_n=\Bigl(\dfrac{1}{i+j-1}\Bigr)_{1\leqslant i,j\leqslant n}$ et
$c_m=\prod_{k=1}^{m-1}k!$ (avec $c_1=1$).
\begin{questions}[resume]
  \item Montrer que $H_n$ est une matrice de Cauchy, puis que
    \[
    \det H_n=\frac{c_n^{4}}{c_{2n}} .
    \]
    Calculer $\det H_2$, $\det H_3$, $\det H_4$ ; en déduire que $H_n$ est
    toujours inversible.
  \item Montrer que le coefficient $(i,j)$ de $H_n$ est
    $\int_0^1x^{i-1}x^{j-1}\,\mathrm{d}x$ : $H_n$ est la matrice de Gram de
    $(1,X,\dots,X^{n-1})$ (chapitre 18).
  \item Calculer $H_3^{-1}$ par la comatrice et constater que ses
    coefficients sont entiers (c'est vrai pour tout $n$, on l'admet).
    Pourquoi la petitesse de $\det H_n$ rend-elle la résolution numérique
    de $H_nX=B$ très instable ?
\end{questions}

\medskip
\begin{center}
\begin{tikzpicture}[scale=0.95]
  \foreach \i in {1,...,6}{\foreach \j in {1,...,6}{%
    \pgfmathsetmacro{\v}{1/(\i+\j-1)}
    \pgfmathsetmacro{\pc}{100*\v}
    \fill[otNavy!\pc!white] (\j-1,-\i+1) rectangle (\j,-\i);
    \ifdim\v pt>0.45pt
      \node[font=\small,text=white] at (\j-0.5,-\i+0.5) {$\frac{1}{\the\numexpr\i+\j-1\relax}$};
    \else
      \node[font=\small,text=otInk] at (\j-0.5,-\i+0.5) {$\frac{1}{\the\numexpr\i+\j-1\relax}$};
    \fi}}
  \draw[otNavy,line width=0.6pt] (0,0) rectangle (6,-6);
\end{tikzpicture}\\[0.5mm]
\legende{La matrice de Hilbert $H_6$ : des coefficients constants le long
des anti-diagonales, qui décroissent lentement — et pourtant
$\det H_6\approx5\cdot10^{-18}$.}
\end{center}
\end{defi}

